% diagram slots are inserted by build_pdf.py
DIAGRAM1
```{=latex}
\begin{figure}[htbp]
\centering
\begin{tikzpicture}[
  font=\small,
  band/.style={draw=line, fill=cream, rounded corners=1.2pt, text width=14.6cm, align=left, inner xsep=9pt, inner ysep=7pt},
  outside/.style={draw=muted, dashed, fill=white, rounded corners=1.2pt, text width=14.6cm, align=left, inner xsep=9pt, inner ysep=7pt},
  node distance=2.6mm
]
  \node[band] (l1) {\textbf{\color{golddeep}渠道层}\enspace 微信小程序：对话、确认卡、适老版字号与语音播放};
  \node[band, below=of l1] (l2) {\textbf{\color{golddeep}接入层}\enspace 登录与会话校验、对话、指令确认与撤销、风险测评、提醒、语音};
  \node[band, below=of l2] (l3) {\textbf{\color{golddeep}智能编排层}\enspace 统一助手、主模型工具调用、语义审核、工具层参数与合规校验};
  \node[band, below=of l3] (l4) {\textbf{\color{golddeep}业务与数据层}\enspace 限额、适当性与反欺诈规则；账务账本；心愿备用金调度；待确认指令};
  \node[outside, below=of l4] (l5) {\textbf{\color{muted}外部能力}\enspace 模型推理服务负责理解与话术；语音合成服务负责播报。余额、限额与入账留在本系统内。};
  \draw[gold, line width=0.7pt, -{Stealth[length=1.7mm]}] (l1.south) -- (l2.north);
  \draw[gold, line width=0.7pt, -{Stealth[length=1.7mm]}] (l2.south) -- (l3.north);
  \draw[gold, line width=0.7pt, -{Stealth[length=1.7mm]}] (l3.south) -- (l4.north);
  \draw[muted, dashed, line width=0.6pt, -{Stealth[length=1.7mm]}] (l4.south) -- (l5.north);
\end{tikzpicture}
\caption{系统逻辑架构}
\end{figure}
```

DIAGRAM2
```{=latex}
\begin{figure}[htbp]
\centering
\begin{tikzpicture}[
  font=\footnotesize,
  zhstep/.style={draw=line, fill=cream, rounded corners=1pt, text width=7.5cm, align=left, inner sep=5pt},
  zhask/.style={draw=gold, fill=white, rounded corners=1pt, text width=7.5cm, align=left, inner sep=5pt},
  zhout/.style={draw=terra, fill=white, rounded corners=1pt, text width=5.5cm, align=left, inner sep=4.5pt},
  node distance=2.3mm and 4mm,
  arr/.style={gold, line width=0.6pt, -{Stealth[length=1.6mm]}}
]
  \node[zhstep] (a) {客户消息进入对话接口，先校验登录态};
  \node[zhask, below=of a] (b) {上一张确认卡仍待确认，且本句是肯定回复？};
  \node[zhout, right=of b] (b2) {重新展示该确认卡，请客户在卡片上输入交易密码};
  \node[zhstep, below=of b] (c) {按话题裁剪工具。主模型自主选择工具，最多 5 步；一旦生成卡片即停止};
  \node[zhask, below=of c] (d) {本句是要取回心愿备用金？};
  \node[zhout, right=of d] (d2) {按列表顺序生成取回确认卡；未指明第几笔则追问};
  \node[zhask, below=of d] (e) {主模型是否已经生成卡片？};
  \node[zhout, right=of e] (e2) {返回最后一张卡片};
  \node[zhask, below=of e] (f) {用途、金额、用钱日是否已由客户说全？};
  \node[zhout, right=of f] (f2) {生成心愿备用金确认卡};
  \node[zhask, below=of f] (g) {语义审核：漏调查看、虚假办结，或正常放行};
  \node[zhstep, below=of g] (h) {视审核结论补发查看卡片、改写尚未执行的话术，或在卡号与金额已齐时补转账确认卡};
  \draw[arr] (a) -- (b);
  \draw[arr] (b) -- (c);
  \draw[arr] (c) -- (d);
  \draw[arr] (d) -- (e);
  \draw[arr] (e) -- (f);
  \draw[arr] (f) -- (g);
  \draw[arr] (g) -- (h);
  \draw[terra, line width=0.5pt, -{Stealth[length=1.5mm]}] (b.east) -- (b2.west);
  \draw[terra, line width=0.5pt, -{Stealth[length=1.5mm]}] (d.east) -- (d2.west);
  \draw[terra, line width=0.5pt, -{Stealth[length=1.5mm]}] (e.east) -- (e2.west);
  \draw[terra, line width=0.5pt, -{Stealth[length=1.5mm]}] (f.east) -- (f2.west);
\end{tikzpicture}
\caption{对话处理流程}
\end{figure}
```

DIAGRAM3
```{=latex}
\begin{figure}[htbp]
\centering
\begin{tikzpicture}[x=1cm,y=0.64cm]
  \node[font=\footnotesize\bfseries, text=golddeep] at (0,0.35) {客户};
  \node[font=\footnotesize\bfseries, text=golddeep] at (4.15,0.35) {小程序};
  \node[font=\footnotesize\bfseries, text=golddeep] at (8.3,0.35) {接入服务};
  \node[font=\footnotesize\bfseries, text=golddeep] at (12.45,0.35) {规则与账务};
  \draw[line, line width=0.6pt] (0,0.05) -- (0,-7.15);
  \draw[line, line width=0.6pt] (4.15,0.05) -- (4.15,-7.15);
  \draw[line, line width=0.6pt] (8.3,0.05) -- (8.3,-7.15);
  \draw[line, line width=0.6pt] (12.45,0.05) -- (12.45,-7.15);
  \draw[-{Stealth[length=1.5mm]}, terra] (0,-0.55) -- (4.15,-0.55);
  \node[font=\scriptsize, above] at (2.08,-0.55) {说明卡号与金额};
  \draw[-{Stealth[length=1.5mm]}, terra] (4.15,-1.35) -- (8.3,-1.35);
  \node[font=\scriptsize, above] at (6.22,-1.35) {提交对话};
  \draw[-{Stealth[length=1.5mm]}, terra] (8.3,-2.15) -- (12.45,-2.15);
  \node[font=\scriptsize, above] at (10.38,-2.15) {预校验限额与卡状态};
  \draw[-{Stealth[length=1.5mm]}, gold] (12.45,-2.95) -- (8.3,-2.95);
  \node[font=\scriptsize, above] at (10.38,-2.95) {可出卡，此时不扣款};
  \draw[-{Stealth[length=1.5mm]}, gold] (8.3,-3.75) -- (4.15,-3.75);
  \node[font=\scriptsize, above] at (6.22,-3.75) {展示确认卡};
  \draw[-{Stealth[length=1.5mm]}, terra] (0,-4.55) -- (4.15,-4.55);
  \node[font=\scriptsize, above] at (2.08,-4.55) {输入交易密码};
  \draw[-{Stealth[length=1.5mm]}, terra] (4.15,-5.35) -- (8.3,-5.35);
  \node[font=\scriptsize, above] at (6.22,-5.35) {提交确认};
  \draw[-{Stealth[length=1.5mm]}, terra] (8.3,-6.15) -- (12.45,-6.15);
  \node[font=\scriptsize, above] at (10.38,-6.15) {复核密码、限额与风险提示};
  \draw[-{Stealth[length=1.5mm]}, gold] (12.45,-6.95) -- (4.15,-6.95);
  \node[font=\scriptsize, above, fill=white, inner sep=1pt] at (8.3,-6.95) {入账后返回回执与最新余额};
\end{tikzpicture}
\caption{资金类业务的确认与入账。收款人为本行客户时，同一处理中贷记对方活期。}
\end{figure}
```
